Our results confirm a systematic increase in the performance gap, which aligns with the observed decline in performance at national championships. This reduction in performance density is particularly evident among men and young athletes, where performance differences have widened significantly. A major reason for this structural change is likely the 11.9\% decrease in German Athletics Association (DLV) membership, exacerbated by a 26.6\% drop in the 27- to 60-year-old age group \citep{LeichtathletikPortalDownloads}. Since this age group contains a large number of coaches, its decline directly impacts technical disciplines that rely heavily on expert, individualized training.  In contrast, women’s athletics shows
more improvements, which is probably due to a ”catch-up
effect” as a result of historically delayed professionalization.

Ultimately, the data reveals a polarization of performance: while professionalization protects the national elite, the second tier is shrinking. This confirms our hypothesis that the lack of internal pressure—true to the principle that "competition promotes performance"—has isolated the German elite and ill-prepared them for the intensity of international competition. The zero medal results of recent years are therefore not anomalies, but the predictable outcome of a shrinking talent pool.

In running competitions, the overall improvement in average performance was significantly influenced by the introduction of carbon fiber shoes in 2016 \citep{willwacherDoesAdvancedFootwear2023}. Furthermore, the COVID-19 pandemic led to uneven development. Endurance athletes may benefited from focused, uninterrupted training blocks \citep{venkatesanRunningPerformanceCOVID192022}. In contrast, athletes in technical disciplines were hindered by the closure of training facilities.

One limitation is the focus on the top 35. We cannot make any statements about the general development in all of amateur sports. The analysis only considers advanced competitive sports. Whether the observed trend continues linearly would need to be investigated with further analyses.

The increase in the performance spread indicates a structural change in organized athletics. 
A future focus should be placed on strengthening the grassroots level. 
Following the principle of "the grassroots supporting the elite," this will allow for more top performances and could bring German athletics closer to the world's elite.
